\section{결론}\label{sec:conclusion}
본 연구는 영상 기반 초기 추정기를 고정한 채, 다중 해상도 특징과 알려진 물리적 치수로 코너의 국소 이동을 학습하는 보정 방법을 평가하였다. 이동 후보의 확률 가중평균과 원본 좌표의 상한을 사용하고, 동일한 자세 계산을 통해 코너 수정과 위치·회전 변화의 관계를 비교하였다.

직사각형 실사 319장에서 세 기반의 코너 중앙 오차와 전체 주석 기준 \pckten 이 개선되었으며, 별도 정사각형 119장에서도 2D 보정 효과가 관찰되었다. 반면 ResNet 위치 중앙값의 사후 구간은 0을 포함했고, 일부 P90과 severe 가림의 위치 오차는 악화되었다. 회전 대응 감독의 추가 효과도 일관되지 않았다. 네 기록 영상 8,910프레임에서는 Base와 N3의 자세 산출률이 같은 98.451\%였고 인접 출력 변화는 지표별로 혼합되었다. 이러한 결과는 주어진 정보·학습·참조 조건에서 국소 보정의 적용 가능성과 경계를 함께 보여준다. 앞으로의 판단에는 정사각형 가림·코너 가시성 검수, 리프터의 사람 참조와 객체 대응, 참조 품질 확인 및 독립 관측에서의 평가가 필요하다.
